% GENERATED FILE. Edit canonical lessons, then run npm run generate:books.
% canonical-source-hash: fnv1a64:138d8cb7133bc3d6
% canonical-lessons: KA-C450-mekyanik, KA-C450-gyarej, KA-C450-bidibhaga, KA-C450-kettuhogu, KA-C450-riperi

\chapter{At the Garage}
\label{ch:ka-at-the-garage}
\hlchaptermodality{450}

\begin{chapteropening}
\textbf{By the end of this chapter:} \emph{I can say a mechanic, a garage, a spare part, to break down and repair work.}

Complete the last lesson of chapter 450: I can say a mechanic, a garage, a spare part, to break down and repair work.
\end{chapteropening}

\section[\texorpdfstring{\kn{ಮೆಕ್ಯಾನಿಕ್} (mekyānik)}{ಮೆಕ್ಯಾನಿಕ್ (mekyānik)}]{\kn{ಮೆಕ್ಯಾನಿಕ್} (mekyānik) — a mechanic}

\label{lesson:KA-C450-mekyanik}



\begin{quote}
Before the new one: say the Kannada for to get a puncture, then the Kannada for a brake.
\end{quote}

\subsection*{You'll want to know: \kn{ಮೆಕ್ಯಾನಿಕ್}}
\textbf{\kn{ಮೆಕ್ಯಾನಿಕ್}} — \emph{mekyānik} — \textquotedblleft{}a mechanic\textquotedblright{}.

\begin{quote}
\kn{ಮೆಕ್ಯಾನಿಕ್} \kn{ಬ್ರೇಕನ್ನು} \kn{ಸರಿಪಡಿಸಿದನು}.
\end{quote}

\begin{quote}
\emph{mekyānik brēkannu saripaḍisidanu.}
\end{quote}

\begin{quote}
The mechanic fixed the brake.
\end{quote}

The plural takes -\kn{ಗಳು}: \kn{ಮೆಕ್ಯಾನಿಕ್ಗಳು}.

\begin{cousinweb}
\textbf{\kn{ಮೆಕ್ಯಾನಿಕ್}} (\emph{mekyānik}) is a loan from English \emph{mechanic}.
\end{cousinweb}

\begin{grammarlens}[title={What you've built}]
One more word for getting around.
\end{grammarlens}

\subsection*{Your turn}
\begin{itemize}
\raggedright
  \item \emph{Say it:} \emph{mekyānik}
  \item \emph{Say it:} \emph{mekyānik}, once more
  \item \emph{Say it:} \emph{mekyānik brēkannu saripaḍisidanu.} — The mechanic fixed the brake.
  \item \emph{Recall:} say the Kannada for to get a puncture, then the Kannada for a brake, then say \emph{mekyānik} again
\end{itemize}

\begin{tcolorbox}[breakable,colback=teal!4,colframe=teal!35!black,title={Before you move on}]
What does \kn{ಮೆಕ್ಯಾನಿಕ್} mean? (\textquotedblleft{}A mechanic\textquotedblright{}.) Say it once more.
\end{tcolorbox}
\section[\texorpdfstring{\kn{ಗ್ಯಾರೇಜ್} (gyārēj)}{ಗ್ಯಾರೇಜ್ (gyārēj)}]{\kn{ಗ್ಯಾರೇಜ್} (gyārēj) — a garage; a workshop}

\label{lesson:KA-C450-gyarej}



\begin{quote}
Before the new one: say the Kannada for a brake, then the Kannada for a mechanic.
\end{quote}

\subsection*{You'll want to know: \kn{ಗ್ಯಾರೇಜ್}}
\textbf{\kn{ಗ್ಯಾರೇಜ್}} — \emph{gyārēj} — \textquotedblleft{}a garage; a workshop\textquotedblright{}.

\begin{quote}
\kn{ನಮ್ಮ} \kn{ಬೀದಿಯ} \kn{ಕೊನೆಯಲ್ಲಿ} \kn{ಒಂದು} \kn{ಗ್ಯಾರೇಜ್} \kn{ಇದೆ}.
\end{quote}

\begin{quote}
\emph{namma bīdiya koneyalli ondu gyārēj ide.}
\end{quote}

\begin{quote}
There is a garage at the end of our street.
\end{quote}

In India mostly a repair workshop for vehicles.

\begin{cousinweb}
\textbf{\kn{ಗ್ಯಾರೇಜ್}} (\emph{gyārēj}) is a loan from English \emph{garage}.
\end{cousinweb}

\begin{grammarlens}[title={What you've built}]
One more word for getting around.
\end{grammarlens}

\subsection*{Your turn}
\begin{itemize}
\raggedright
  \item \emph{Say it:} \emph{gyārēj}
  \item \emph{Say it:} \emph{gyārēj}, once more
  \item \emph{Say it:} \emph{namma bīdiya koneyalli ondu gyārēj ide.} — There is a garage at the end of our street.
  \item \emph{Recall:} say the Kannada for a brake, then the Kannada for a mechanic, then say \emph{gyārēj} again
\end{itemize}

\begin{tcolorbox}[breakable,colback=teal!4,colframe=teal!35!black,title={Before you move on}]
What does \kn{ಗ್ಯಾರೇಜ್} mean? (\textquotedblleft{}A garage; a workshop\textquotedblright{}.) Say it once more.
\end{tcolorbox}
\section[\texorpdfstring{\kn{ಬಿಡಿಭಾಗ} (biḍibhāga)}{ಬಿಡಿಭಾಗ (biḍibhāga)}]{\kn{ಬಿಡಿಭಾಗ} (biḍibhāga) — a spare part}

\label{lesson:KA-C450-bidibhaga}



\begin{quote}
Before the new one: say the Kannada for a mechanic, then the Kannada for a garage.
\end{quote}

\subsection*{You'll want to know: \kn{ಬಿಡಿಭಾಗ}}
\textbf{\kn{ಬಿಡಿಭಾಗ}} — \emph{biḍibhāga} — \textquotedblleft{}a spare part\textquotedblright{}.

\begin{quote}
\kn{ಹಳೆಯ} \kn{ಕಾರಿನ} \kn{ಬಿಡಿಭಾಗಗಳು} \kn{ಸಿಗುವುದು} \kn{ಕಷ್ಟ}.
\end{quote}

\begin{quote}
\emph{haḷeya kārina biḍibhāgagaḷu siguvudu kaṣṭa.}
\end{quote}

\begin{quote}
Spare parts for an old car are hard to find.
\end{quote}

Plural: \kn{ಬಿಡಿಭಾಗಗಳು}, \emph{biḍibhāgagaḷu}.

\kn{ಬಿಡಿ}, loose or separate, plus \kn{ಭಾಗ}, a part.

\begin{cousinweb}
\textbf{\kn{ಬಿಡಿಭಾಗ}} (\emph{biḍibhāga}) joins native \emph{biḍi}, \textquotedblleft{}loose, separate\textquotedblright{}, from \emph{biḍu}, \textquotedblleft{}to let go\textquotedblright{} (DEDR 5393), and Sanskrit \emph{bhāga}, \textquotedblleft{}a portion, a part\textquotedblright{} (CDIAL 9430).
\end{cousinweb}

\begin{grammarlens}[title={What you've built}]
One more word for getting around.
\end{grammarlens}

\subsection*{Your turn}
\begin{itemize}
\raggedright
  \item \emph{Say it:} \emph{biḍibhāga}
  \item \emph{Say it:} \emph{biḍibhāga}, once more
  \item \emph{Say it:} \emph{haḷeya kārina biḍibhāgagaḷu siguvudu kaṣṭa.} — Spare parts for an old car are hard to find.
  \item \emph{Recall:} say the Kannada for a mechanic, then the Kannada for a garage, then say \emph{biḍibhāga} again
\end{itemize}

\begin{tcolorbox}[breakable,colback=teal!4,colframe=teal!35!black,title={Before you move on}]
What does \kn{ಬಿಡಿಭಾಗ} mean? (\textquotedblleft{}A spare part\textquotedblright{}.) Say it once more.
\end{tcolorbox}
\section[\texorpdfstring{\kn{ಕೆಟ್ಟುಹೋಗು} (keṭṭuhōgu)}{ಕೆಟ್ಟುಹೋಗು (keṭṭuhōgu)}]{\kn{ಕೆಟ್ಟುಹೋಗು} (keṭṭuhōgu) — to break down; to go bad}

\label{lesson:KA-C450-kettuhogu}



\begin{quote}
Before the new one: say the Kannada for a garage, then the Kannada for a spare part.
\end{quote}

\subsection*{You'll want to know: \kn{ಕೆಟ್ಟುಹೋಗು}}
\textbf{\kn{ಕೆಟ್ಟುಹೋಗು}} — \emph{keṭṭuhōgu} — \textquotedblleft{}to break down; to go bad\textquotedblright{}.

\begin{quote}
\kn{ಹೆದ್ದಾರಿಯ} \kn{ಮಧ್ಯದಲ್ಲಿ} \kn{ನಮ್ಮ} \kn{ಬಸ್} \kn{ಕೆಟ್ಟುಹೋಯಿತು}.
\end{quote}

\begin{quote}
\emph{heddāriya madhyadalli namma bas keṭṭuhōyitu.}
\end{quote}

\begin{quote}
Our bus broke down in the middle of the highway.
\end{quote}

Past (it …): \kn{ಕೆಟ್ಟುಹೋಯಿತು}, \emph{keṭṭuhōyitu}.

Of food it means to spoil: \kn{ಹಾಲು} \kn{ಕೆಟ್ಟುಹೋಯಿತು}, the milk went off.

\begin{cousinweb}
\textbf{\kn{ಕೆಟ್ಟುಹೋಗು}} (\emph{keṭṭuhōgu}) is native Dravidian: \emph{keṭṭu}, the past stem of \emph{keḍu}, \textquotedblleft{}to be ruined, to spoil\textquotedblright{} (DEDR 1942), with \kn{ಹೋಗು}, to go.
\end{cousinweb}

\begin{grammarlens}[title={What you've built}]
One more word for getting around.
\end{grammarlens}

\subsection*{Your turn}
\begin{itemize}
\raggedright
  \item \emph{Say it:} \emph{keṭṭuhōgu}
  \item \emph{Say it:} \emph{keṭṭuhōgu}, once more
  \item \emph{Say it:} \emph{heddāriya madhyadalli namma bas keṭṭuhōyitu.} — Our bus broke down in the middle of the highway.
  \item \emph{Recall:} say the Kannada for a garage, then the Kannada for a spare part, then say \emph{keṭṭuhōgu} again
\end{itemize}

\begin{tcolorbox}[breakable,colback=teal!4,colframe=teal!35!black,title={Before you move on}]
What does \kn{ಕೆಟ್ಟುಹೋಗು} mean? (\textquotedblleft{}To break down; to go bad\textquotedblright{}.) Say it once more.
\end{tcolorbox}
\section[\texorpdfstring{\kn{ರಿಪೇರಿ} (ripēri)}{ರಿಪೇರಿ (ripēri)}]{\kn{ರಿಪೇರಿ} (ripēri) — repair work; a repair}

\label{lesson:KA-C450-riperi}



\begin{quote}
Before the new one: say the Kannada for a spare part, then the Kannada for to break down.
\end{quote}

\subsection*{You'll want to know: \kn{ರಿಪೇರಿ}}
\textbf{\kn{ರಿಪೇರಿ}} — \emph{ripēri} — \textquotedblleft{}repair work; a repair\textquotedblright{}.

\begin{quote}
\kn{ಗ್ಯಾರೇಜ್ನಲ್ಲಿ} \kn{ಕಾರಿನ} \kn{ರಿಪೇರಿ} \kn{ಎರಡು} \kn{ದಿನ} \kn{ನಡೆಯಿತು}.
\end{quote}

\begin{quote}
\emph{gyārējnalli kārina ripēri eraḍu dina naḍeyitu.}
\end{quote}

\begin{quote}
The repair of the car at the garage took two days.
\end{quote}

\kn{ರಿಪೇರಿ} \kn{ಮಾಡು}, to repair, sits beside the Kannada \kn{ಸರಿಪಡಿಸು}.

\begin{cousinweb}
\textbf{\kn{ರಿಪೇರಿ}} (\emph{ripēri}) is a loan from English \emph{repair}, reshaped with a final -\emph{i}.
\end{cousinweb}

\begin{grammarlens}[title={What you've built}]
One more word for getting around.
\end{grammarlens}

\subsection*{Your turn}
\begin{itemize}
\raggedright
  \item \emph{Say it:} \emph{ripēri}
  \item \emph{Say it:} \emph{ripēri}, once more
  \item \emph{Say it:} \emph{gyārējnalli kārina ripēri eraḍu dina naḍeyitu.} — The repair of the car at the garage took two days.
  \item \emph{Recall:} say the Kannada for a spare part, then the Kannada for to break down, then say \emph{ripēri} again
\end{itemize}

\begin{tcolorbox}[breakable,colback=teal!4,colframe=teal!35!black,title={Before you move on}]
What does \kn{ರಿಪೇರಿ} mean? (\textquotedblleft{}Repair work; a repair\textquotedblright{}.) Say it once more.
\end{tcolorbox}
